\capitulo{5}{Aspectos relevantes del desarrollo del proyecto}

Este capítulo recoge los aspectos más relevantes del desarrollo práctico del proyecto, poniendo el foco en aquellas decisiones de análisis, diseño e implementación que han condicionado de forma significativa la solución final. Su finalidad no es reproducir de manera literal fragmentos de código o elementos ya presentes en los anexos, sino justificar los principales caminos de solución adoptados.

\section{Modelo híbrido de persistencia}

Uno de los aspectos más relevantes del desarrollo ha sido el integrar dos enfoques distintos de persistencia dentro de un mismo sistema. Por un lado, la parte de la base de datos asociada a la gestión académica y organización de la escuela que se ha realizado a través de un \textit{script} SQL cuya base o prototipo fue realizado al principio del proyecto, previo a elegir \textit{frameworks} o lenguajes de programación. Y por otro lado, tenemos la gestión de usuarios, roles y permisos se implementó con los mecanismos que proporciona Django.

Esta decisión permitió reutilizar la estructura de datos preexistente, evitando redefinir las entidades ya consolidadas en el \textit{script}, y a su vez aprovechar las facilidades que ofrece el \textit{framework} para la autenticación, autorización y administración de accesos. No obstante, hubo que integrar las entidades, relaciones y reglas de acceso para aprovechas las ventajas del ORM que ofrece Django. Para ello se volcaron los modelos a Django mediante el comando \textit{inspectdb}.

Desde el punto de vista práctico, este modelo híbrido permitió combinar una base de datos académica ya definida con un sistema de seguridad y gestión de usuarios más flexible, delegando la lógica de acceso a la aplicación en un servicio externo, y manteniendo en la base de datos propia la asignación de roles y permisos a cada usuario.

\section{Integración con la autenticación institucional}

Otro aspecto destacado del desarrollo fue la decisión de delegar la autenticación de usuarios en la API institucional de la universidad. De este modo el sistema no valida directamente las credenciales contra la base de datos propia, sino que utiliza las credenciales oficiales proporcionadas por la institución, apoyándose así en un mecanismo externo ya consolidado y en funcionamiento.

Esta decisión fue tomada debido a que la aplicación va a ser utilizada y gestionada por personas relacionadas a la Universidad de Burgos, de forma que así se controla los usuarios que pueden acceder a la página y evitando así tener que hacer un registro de usuario, provocando que el usuario tenga que crear otra contraseña nueva para tareas relacionadas con la institución. Por lo que se llegó a la conclusión de que este paso no era necesario.

Para la integración se tomó de ejemplo su implementación en otro proyecto programado en Java, que también delegaba la autenticación a la API del moodle. Con este proyecto como base se siguieron los pasos necesarios para su correcta integración.

De esta forma, se dividieron las tareas, delegando así el proceso de autenticación de usuarios y validación de credenciales a la API del moodle de la Universidad e Burgos. Una vez el usuario se autentica de forma correcta, se comprueba en la base de datos propia la existencia de una entrada para el usuario logueado para poder obtener el rol de este, dándole acceso a las operaciones permitidas según los permisos asignados que tenga su rol. Si no se encuentra una entrada en la base de datos para el usuario se decidió crear una entrada en base de datos de manera automática pero asignándole el rol de PDI, es decir, el rol con menos privilegios asignados. De forma que no pueda realizar operaciones críticas en el sistema, y manteniendo así la seguridad e integridad del mismo.

Esto ha permitido combinar un acceso institucional unificado a una gestión interna de permisos ajustada a las necesidades reales del sistema.

\section{Gestión del calendario académico}

Otro de los aspectos relevantes del desarrollo fue la necesidad de incorporar al sistema la gestión del calendario académico. Durante el análisis del problema se observó que la simple definición de reservas periódicas no era suficiente para poder representar las clases, ya que la ocupación efectiva de las aulas depende también del tipo de día académico sobre el que se proyecta cada actividad.

En el contexto universitario, no todos los días del curso tienen el mismo comportamiento. Existen jornadas de docencia lectiva, periodos de exámenes, días festivos, cambios de docencia y días dedicado exclusivamente a Trabajo de Fin de Grado y Trabajos de Fin de Máster. Por este motivo, se consideró necesario añadir esta información dentro de la propia aplicación, evitando depender de una interpretación estática de los horarios.

Como respuesta a esta necesidad, se introdujo la entidad curso académico, estructurado a su vez en semestres, y asociado a una clasificación de los días y sus tipos. Con esta organización se buscó representar de una forma más precisa y adaptada la planificación temporal al contexto en el que se desarrollan las actividades docentes y permite que el sistema identifique cuándo debe tenerse en cuenta o añadirse una reserva periódica y cuando no.

La gestión relacionada a los tipos de días aporta varias ventajas. En primer lugar, permite que la consulta de ocupación de aulas garantice la realidad académica. En segundo lugar, mejora la generación de reservas periódicas, de forma que estas se adaptan al calendario académico efectivo en lugar de repetirse de manera uniforme sobre todos los días de la semana. Por último, introduce flexibilidad dentro del sistema, debido a que permite gestionar desde la interfaz en situaciones excepcionales como festivos, periodos de exámenes o modificaciones puntuales de la docencia.

Desde el punto de vista del diseño, esta decisión supone que el sistema no solo gestiona reservas y horarios, sino también se encarga del marco temporal académico que condiciona la validez de los puntos anteriores. En consecuencia, la aplicación pasa a tratar no únicamente qué actividad tiene lugar en el aula, sino también en qué contexto académico se produce, lo que refuerza su adecuación al entorno universitario y mejora la calidad de la información que ofrece.

\section{Parser de horarios}
Durante el desarrollo del proyecto, se decidió realizar un analizador léxico para la carga de los horarios de los grados en el sistema, los cuales se encuentran ahora mismo en una hoja de cálculo. Esta decisión fue tomada debido a que para poder modificar la franja horario en la que se imparte una clase tiene que estar cargada en base de datos, y habría que cargar de manera manual las reservas en el sistema. Lo que no es viable ya que con la extracción de los datos de la hoja de cálculo, el número de clases o reservas extraídas supera las 1000 reservas periódicas.

Para comenzar se analizó en profundidad la estructura que presentaban las clases teóricas y prácticas de la hoja de cálculo que se utiliza actualmente para la gestión de los horarios. Tras el análisis se supuso que las clases teóricas podían ser encontradas con tres estructuras diferentes, y las clases prácticas con otras tres estructuras, con una estructura de ellas común entre los dos tipos de clases. También se observó que las clases teóricas tenían un color y las prácticas otro a nivel general incluyendo las clases en inglés del grado de Organización Industrial. También se analizaron los datos que era importante extraer para poder realizar las cargas en base de datos de las reservas periódicas, entre los que se encuentran, la asignatura, su código (identificador de la asignatura en base de datos), el grupo, el aula, el día de la semana y la franja horaria en la que se imparte la clase. 

Durante el desarrollo de este \textit{parser}, se detectó que la hoja de cálculo inicial presentaba una estructura poco homogénea. Aunque a nivel general, sí que seguía una estructura, se encontraron datos introducidos sin mecanismos de validación consistentes, y en ocasiones sin un criterio uniforme. Por consecuencia, se realizó una fase de depuración, normalización y unificación de la información, que supuso una cantidad de tiempo considerable a lo largo del desarrollo para que pudiera ser integrada de manera correcta en el sistema.
Para poder comprender el tiempo invertido cabe destacar que los valores cargados en los horarios se introducían o modificaban a través de validadores de datos, en este caso listas o desplegables. Con lo que cada vez que se modificaba el valor de uno de estos datos en la hoja correspondiente, había que ir a las hojas de los horarios respectivas en las que se encontraba el dato y volver a seleccionarlo para que se actualizada, esto para cada una de las celdas que contenía ese dato.

Entre los problemas de unificación encontrados están los siguientes:

\begin{enumerate}
	\item Inconsistencia en el formato de colores. Aunque visualmente había colores que podían parecer similares, la tonalidad o código era distinto para representar una misma categoría. El grupo teórico de inglés tenía el mismo color designado que el grupo teórico para el segundo grupo de teoría.
	\begin{itemize}
		\item Solución aplicada: Se decidió qué color y tonalidad iba a tener cada categoría encontrada, y 
	\end{itemize}
	\item Introducción manual de datos. Se encontraron datos que en vez de haber sido añadidos mediante el validador se habían introducido a mano en la celda, en vez de ser añadidos en la lista de celdas de la que el validador extraía los datos.
	\begin{itemize}
		\item Solución aplicada: Se analizó y detectó que datos habían sido introducidos de manera manual, se añadió su valor a la lista del validador, y se eligió mediante este.
	\end{itemize}
	\item Inconsistencia en la nomenclatura de los laboratorios. Para referirse a los laboratorios, se utilizaban diferentes abreviaturas (Lab., Lab, Lb., L.). La última de estas abreviaturas provocaba que alguna asignatura la identificase como aula por acabar con <<l>>. Por ejemplo, APL. BASES DE DATOS (Aplicación de base de datos).
	\begin{itemize}
		\item Solución aplicada: Se unificó la abreviatura a Lab. o Lab para todos los laboratorios. Y posteriormente, se tuvo que actualizar en cada celda en la que se encontraba uno de los laboratorios modificados. 
	\end{itemize}
	\item Inconsistencia en la separación de aulas combinadas. Las combinaciones de aulas eran separadas por diferentes caracteres especiales (<<,>>, <</>>, <<->>)
	\begin{itemize}
		\item Solución aplicada: Se unificó el separador de las aulas a <<\textit{/}>> en el validador, y se actualizó en cada celda que aparecía la combinación de todos los horarios para que se aplicaran los cambios.
	\end{itemize}
	\item Inconsistencia en las abreviaturas de las asignaturas. Se encontró que las abreviaturas de las asignaturas no eran homogéneas. La misma asignatura podía encontrarse en el mismo horario con tilde y sin ella, o directamente la abreviatura fuese totalmente diferente. Por ejemplo, se podía encontrar: AMP. CALC., AMP. CÁLCULO o AMP. CALCULO
	\begin{itemize}
		\item Solución aplicada: Unificar la abreviatura de las asignaturas en la hoja de Asignaturas, e ir actualizando en el horario del grado y semestre correspondiente.
	\end{itemize}
	\item Valores con fuente blanca. Se encontraron celdas en las que el fondo era blanco y la fuente también, es decir, estos valores existían pero no se veían.
	\begin{itemize}
		\item Solución aplicada: Se detectaron que valores se imprimían que no entraban dentro de las estructuras predefinidas con colores, y se puso el color de la fuente a negro y se eliminaron los valores.
	\end{itemize}
	\item Filas ocultas. Se encontraron filas ocultas que estaban dando error porque no tenían los nombre unificados, y tras tiempo analizando bien la situación resultó que no se veían y estaban ocultas.
	\begin{itemize}
		\item Solución aplicada: Se expandieron todas las filas en las que se producía salto y no eran contiguas, y se eliminaron.
	\end{itemize}
	\item Clases sin aula asignada. Algunas clases no tienen definida el aula por no conocerse si finalmente van a ser impartidas o no, corresponden a asignaturas optativas. Pero para poder realizar la reserva, es obligatorio introducir un aula.
	\begin{itemize}
		\item Solución aplicada: Se implementó en las distintas estructuras del analizador léxico que si no se encontraba un valor de tipo Aula en la celda correspondiente se le iba a asignar el valor Aula 0. Y en base de datos se añadió un Aula 0 que no existe, solamente es provisional para poder realizar la reserva.
	\end{itemize}
\end{enumerate}

Se ha adjuntado a continuación una imagen~\ref{fig:problemas-unificaciones}, en la que vamos a poder observar y analizar varios de los problemas que acaban de ser expuestos.

\imagen{img/InconsistenciasExcel.png}{Imagen apreciativa de los problemas de uniformidad de la hoja de cálculo}{fig:problemas-unificaciones}

En color azul, se puede observar como en la clase teórica aparece la asignatura sin tilde y separada por un guión del I, en cambio, en la clase práctica se puede observar como aparece la asignatura sin tilde y sin el guión de separación.
De color verde se ve un caso similar, entre dos clases teóricas una de ellas sin tilde y abreviada, y la otra por el contrario con el nombre completo y la tilde asignada, algo similar al caso de color rojo o marrón.
De color amarillo, se observa como en uno de los casos el separador es un guión y en el otro una barra. Además, fijándose en pequeños detalles que pueden pasar desapercibidos a simple vista se puede llegar a apreciar como en la primera clase Inf 1B referido al aula se encuentra separado, y en la impartida el jueves se encuentra todo seguido Inf1B.

En la siguiente imagen~\ref{fig:problemas-colores}, se van a apreciar como había diferentes tonos de azul asignados para las clases prácticas del mismo grupo, en este caso, el grupo 1.

\imagen{img/ColoresClasesExcel.png}{Imagen apreciativa de los diferentes tonos de color asignados a un mismo tipo de clase y grupo}{fig:problemas-colores}


Y por último, se van a adjuntar unas imágenes~\ref{fig:problemas-estructuras-1,fig:problemas-estructuras-2,fig:problemas-estructuras-3}, en las que se van a poder ver las diferentes estructuras que se han tenido que tener en cuenta para poder extraer de manera correcta la información de las clases.

\imagen{img/EstructurasClases.png}{Imagen con dos tipos de estructuras para los grupos teóricos y los prácticos}{fig:problemas-colores}

En esta primera imagen se observan dos tipos de clases prácticas, ambas con el texto rotado 90º. Una con el grupo arriba de la asignatura, y la otra con el grupo a la derecha. Además se aprecian las clases teóricas con el texto sin rotación y con texto rotado 90º al igual que las prácticas. Para poder encontrarlo mejor, el color morado y carne se refieren a clases teóricas, y los colores verde y azul a clases prácticas.

\imagen{img/EstructurasClases2.png}{Imagen para ver estructuras de grupos teóricos cuando hay 2}{fig:problemas-estructuras-2}

En esta segunda imagen se puede apreciar que cuando hay dos grupos teóricos, en vez de aparecer el aula debajo de la asignatura aparece a la derecha, además de que también se debe guardar el número del grupo que aparece al final del bloque.

\imagen{img/EstructurasClases3.png}{Imagen donde se aprecia otra estructura para los grupos prácticos}{fig:problemas-estructuras-3}

En la última imagen, se observa el último tipo de estructura para los grupos prácticos, similar a una de la de los grupos teóricos, con el nombre de la asignatura y el aula debajo, sin apreciar el número del grupo que por defecto se asignará el 101.

Tras terminar la fase preparación y unificación de la hoja de cálculo, y finalizar de forma correcta la implementación del \textit{parser}, para poder cargar las reservas en base de datos, hubo que cargar previamente todas las asignaturas y aulas con sus respectivos campos en base de datos.

Para poner en contexto la complejidad y el esfuerzo requeridos por el desarrollo del analizador léxico y por la fase previa de unificación de la hoja de cálculo, conviene señalar que este trabajo se prolongó desde el 25 de febrero hasta el 13 de mayo, abarcando un total de 11 semanas. En términos de planificación, este proceso ocupó desde el sprint 13 hasta el sprint 21, es decir, 9 sprints en total, debido a los distintos problemas e inconsistencias que fueron surgiendo durante su desarrollo.


\section{Roles y permisos}

Dado que la aplicación está dirigida a distintos tipos de usuarios con responsabilidades diferenciadas, fue necesario diseñar un modelo de control de acceso basado en roles y permisos. Este aspecto no solamente afecta a la seguridad, sino también a la organización funcional del sistema, ya que determina qué vistas puede consultar cada usuario y qué operaciones puede ejecutar.

La aplicación distingue entre los siguientes roles: administrador, gestor editor, gestor no editor y personal docente e investigador, cada uno de ellos con un conjunto concreto de capacidades.

Se van a distinguir los permisos de los que dispone cada uno de los roles.
\begin{itemize}
	\item Administrador (admin):
	\begin{itemize}
		\item Acceder al apartado de administración
		\item Administrar usuarios, roles y permisos
		\item Cargar un nuevo horario en base de datos
		\item Todas las acciones que pueden realizar los siguientes roles
	\end{itemize}
	\item Gestor editor:
	\begin{itemize}
		\item Modificar clases de los horarios y crear nuevas
		\item Crear y modificar nuevos cursos académicos
		\item Todas las acciones que realizan los siguientes roles
	\end{itemize}
	\item Gestor no editor:
	\begin{itemize}
		\item Consultar horarios
		\item Consultar ocupación de aulas
		\item Consultar reservas puntuales
		\item Gestión de reservas puntuales: crear, editar (aceptar, rechazar, modificar datos) y eliminar.
	\end{itemize}
	\item Personal Docente e Investigador (PDI):
	\begin{itemize}
		\item Solicitar una nueva reserva
		\item Consultar el estado de sus reservas
		\item Editar mientras que el estado de su reserva sea Pendiente
	\end{itemize}
\end{itemize}

Dvidiéndose así, de forma clara las operaciones que puede realizar cada uno de los roles y las funcionalidades a las que puede acceder, adaptándose así a los diferentes tipos de perfiles que van a hacer uso de esta herramienta.



